% !TeX root = ../libro.tex
% !TeX encoding = utf8

\chapter{La aparición de Turing}
\dictum[Bruce Schneier]{La criptografía convierte los datos en información, y la información en poder.}

\section{El contexto}
Alan Turing aparecía en escena de una manera tímida, pero no por ello menos predispuesta. En 1931 entraba en su primera clase de la licenciatura de matemáticas en la Universidad de Cambridge, 
y a partir de ese momento todo fue un proceso rodado. Durante sus estudios, pudo conocer al que sería su mentor durante los próximos años, Max Neumann, el cual le dirigiría en los años venideros 
hacia el campo de la lógica matemática y la teoría de la computación\footnote{Max Newman no sólo sería su mentor en los años venideros, sino que también sería su jefe durante la Segunda Gerra Mundial,
en el centro de investigación criptográfica de Bletchley Park, donde juntos establecieron las bases teóricas del funcionamiento de ``The Bomb''. Véase \cite{copeland2006colossus}.}.

Fue en sus primeros años de estudio durante los que se vivió la caída del programa de Hilbert por parte de Gödel, y fue también este evento lo que motivó a Turing a enfrentarse a este grandioso 
reto que era la decidibilidad.

Pero lo cierto era que la decidibilidad no fue más que una simple aplicación de algo mucho más grande que había diseñado muy detalladamente durante varios años.

\section{El comienzo de la máquina}

Turing presentó durante toda su vida un gran interés en la mente humana y en la posibilidad de que la mente pudiera ser considerada como una máquina de computación. Este joven matemático creía firmemente que 
la mente humana era capaz de realizar cálculos y procesos similares a los que se llevaban a cabo en las máquinas, y que estas capacidades podrían ser estudiadas y entendidas a través 
de la teoría de la computación que había aprendido con su mentor.

Este gran visionario del siglo XX argumentó en más de una ocasión que creía firmemente que cualquier problema que pudiera ser resuelto por un ser humano podría ser resuelto por una máquina de computación, siempre y cuando 
se pudiera describir un algoritmo que resolviera el problema. En otros términos, consideraba que la mente humana y las máquinas podían ser conceptualmente similares en términos de su capacidad de realizar cálculos y 
procesos\footnote{Sin ser consciente de ello, este gran matemático tenía en mente el concepto de inteligencia artificial del que tanto hablamos hoy en día, pero 80 años atrás. 
Tanta fue su ambición que diseñó un test para comprobar cómo de inteligente podía ser una máquina. Véase \cite{turing1950computing}.}.

Turing estuvo varios años buscando una forma de describir estas máquinas de computación y establecer reglas precisas para su funcionamiento. Fue así como, tras varios intentos fallidos,
llegó a la creación de la máquina de Turing, que se ha convertido en un concepto fundamental en el campo de la informática y la teoría de la computación.

Esta máquina teórica puede leer, escribir y borrar símbolos en la cinta, y cambiar de estado según las reglas predefinidas\footnote{Más adelante trataremos en profundidad acerca de las Máquinas de Turing y el concepto de
el concepto de máquina universal.}. La idea detrás de la máquina de Turing es que cualquier cálculo o proceso 
algorítmico se puede simular utilizando este modelo teórico. 

Y a pesar de que fue un gran hito para Turing, no dejaba de ser eso, un modelo teórico. Él quería llevarlo mucho más allá: quería darle vida.

\section{Las máquinas cobran vida}

En su afán por describir físicamente estas máquinas, Turing llegó a construir varios prototipos de computadores que nunca llegaron a buen puerto. Sin embargo, no pasó igual con dos proyectos desarrollados 
durante la Segunda Guerra Mundial, y que por supuesto cambiaron el rumbo de la historia: ``The Bomb'' y ``Colossus''.

Ambas fueron dos máquinas de descifrado de mensajes encriptados del bando alemán por parte de los aliados\footnote{Véase \cite{davis2000universal}.}. Cada una de ellas estaba especializada en una máquina de cifrado alemán: Enigma y Lorenz. Es cierto
que en el desarrollo de ``The Bomb'', Turing tuvo un papel más en la sombra en el desarrollo de la máquina, aportando sus conocimientos no tanto en el diseño, sino más bien en la parte 
criptográfica\footnote{Cabe destacar que, como en todos los grandes proyectos, fueron muchas las personas involucradas en este (aun más si cabe debido al calibre del problema). Véase \cite{sale2015enigma}.}, la cual se fundamentaba en sus conocimientos de lógica y criptografía estudiados durante tantos años.

Sin embargo, con ``Colossus'' fue diferente. La máquina de cifrado Lorenz era mucho más potente que Enigma\footnote{Para profundizar en el desarrollo de la resolución del problema para el descifrado de la máquina Enigma, véase \cite{hodges1983alan}}, y se requería para poder abatirla de toda la élite británica de la época. Turing, aunando sus conocimientos
en criptografía y sus conocimientos electromecánicos previos, se dispuso a la cabeza del equipo de diseño y construcción de la máquina. 

Pero este gran matemático no tenía como único objetivo descifrar mensajes: 
su gran motivación era la máquina en sí misma. Turing veía la oportunidad de aplicar los principios de su Máquina de Turing para construir una máquina especializada, el Colossus, que pudiera realizar tareas de 
criptoanálisis más rápidamente y de manera más eficiente.

En cierto sentido, Turing veía al Colossus como un paso hacia su visión más amplia de una máquina universal de computación\footnote{Realmente, a pesar del patriotismo que presentaba Turing, su principal motivación era llevar a
cabo una aproximación lo más real posible de su máquina universal. Véase \cite{hodges1983alan}}. Aunque el Colossus era una máquina especializada para el descifrado de mensajes de Lorenz, 
compartía algunos principios clave con la Máquina de Turing, como la capacidad de procesar información mediante la manipulación de señales eléctricas y el uso de técnicas de programación para ajustar sus operaciones.

\section{Sentando las bases del futuro}

A pesar de que Colossus no fuera una máquina universal en el sentido más amplio, Turing pudo aplicar sus conocimientos y perspectiva teórica en su diseño y desarrollo. Su experiencia con la Máquina de Turing y su 
deseo de construir máquinas que pudieran realizar cálculos complejos influyeron en su enfoque y en la búsqueda de soluciones innovadoras durante la construcción del Colossus.

Aunque no llegó a ser un computador como tal, a nivel arquitectónico llegó a ser uno de los mayores hitos hasta el momento en el campo de la criptografía, sentando las bases del funcionamiento de los computadores actuales. Y no sólo
encaminó el avance del campo que hoy en día conocemos como informática, sino que con ella pudo salvar miles de vidas durante la Segunda Guerra Mundial por parte del bando aliado, contribuyendo a que la duración de
la guerra se redujera significativamente.


\chapter{Máquina de Turing}

\section{El concepto de Máquina de Turing}

A nivel conceptual, la máquina de Turing es un modelo matemático que pretende resolver un problema que se le introduce como entrada, con la intención de proporcionar una salida pasado un cierto número de pasos.
Es decir, un modelo capaz de simular cualquier algoritmo de computación que se le puedan pasar como instrucciones a la máquina.

La máquina de Turing consta de una cinta infinita dividida en casillas, una cabezal de lectura/escritura que puede moverse a lo largo de la cinta, y un conjunto de reglas que determinan cómo la máquina puede 
cambiar el contenido de las casillas y mover el cabezal.

En su forma original, Turing define los siguientes elementos que serán los componentes de la máquina:

\begin{enumerate}
    \item Cinta: La cinta es una estructura lineal dividida en casillas, donde cada casilla puede contener un símbolo de un alfabeto finito. La cinta es infinita en ambas direcciones, pero en la 
            práctica sólo se usa en un segmento infinito a la derecha.
    \item Cabezal de lectura/escritura: El cabezal de la máquina de Turing puede leer el símbolo actual en la casilla en la que se encuentra y puede escribir un nuevo símbolo en esa casilla. También 
            puede moverse a la casilla adyacente izquierda o derecha en la cinta.
    \item Conjunto de estados: La máquina de Turing tiene un conjunto finito de estados internos que determinan su comportamiento. En cada paso, la máquina de Turing se encuentra en un estado particular.
    \item Reglas de transición: Las reglas de transición especifican cómo la máquina de Turing debe cambiar de estado y cómo debe modificar los símbolos en la cinta en función del estado actual y del símbolo 
            leído por la cabeza de lectura/escritura.
\end{enumerate}

Este concepto revolucionario cambió por completo el curso de la historia, dando paso a lo que hoy conocemos como computadores, pudiendo ver claras similitudes con los procesos que realizan los ordenadores 
de hoy en día. 
\begin{enumerate}
    \item El análogo actual de la cinta, podríamos pensarlo como la memoria de los computadores. Almacena la información de manera ordenada, para un fácil acceso con las órdenes oportunas.
    \item El cabezal de la máquina es el actual procesador del ordenador, el cual realiza operaciones y manipula la información almacenada en la memoria.
    \item El conjunto de estados se complementa junto con el cabezal para componer en conjunto el procesador del ordenador, almacenando los estados de las operaciones que se realizan por el cabezal.
    \item Las reglas de transición se pueden ver como las instrucciones que contiene cada uno de los programas para que pueda ser ejecutado de forma correcta. 
\end{enumerate}

Sin embargo, esto no deja de ser un simple concepto de lo que conocemos como máquina de Turing. En realidad, y formalmente hablando, se trata de un sistema matemático un poco más complejo de describir.

\section{La formalización}

Una máquina de Turing, visto desde un punto de vista formal, es un modelo teórico capaz de realizar cálculos y procesos algorítmicos. Este modelo se define como una 7-tupla $M = (Q, \Sigma, \Gamma, \delta, q_0, q_A, q_R)$,
donde:

\begin{enumerate}
        \item $Q$ es un conjunto finito de estados.
        \item $\Sigma$ es el conjunto finito de símbolos de entrada, también conocido como alfabeto de la máquina\footnote{Denotar que el espacio en blanco ($\sqcup$) no se considera un elemento del alfabeto, ya que es el
                        es el resultado de borrar cualquier elemento del propio alfabeto en la cinta.}.
        \item $\Gamma$ es el conjunto de símbolos presentes en la cinta. Notemos que $\Sigma \in \Gamma$\footnote{Notemos que esto ocurre debido al espacio en blanco ($\sqcup$), ya que $\sqcup \notin \Sigma$, pero 
                        $\sqcup \in \Gamma$, ya que en la cinta, al ser infinita, los esapcios no rellenos por caracteres del alfabeto quedan rellenos por espacios en blanco.}.
        \item $\delta$ es una función de transición que especifica cómo debe comportarse la máquina. Para cada estado y símbolo en la cinta, esta función nos indica cuál será el siguiente estado, el símbolo que 
                        se debe escribir en la cinta y hacia dónde se debe desplazar el cabezal, ya sea izquierda o derecha. A efectos generales, esta función es el corazón de la máquina, ya que nos indica cuál será 
                        su siguiente movimiento. La función se define de la siguiente manera:
                        $$\delta: Q \times \Gamma \longrightarrow Q \times \Gamma \times \{\text{L,R}\}\footnote{La función tomará el valor L o R dependiendo si el cabezal, tras realizar la acción que determine la función,
                        debe moverse a izquierda (L) o a la derecha (R).}$$
                        Es decir, para una entrada $(q,n)$, donde $q$ es el estado de la máquina y $n$ es el símbolo que se 
                        encuentra en ese momento bajo el cabezal, la función devuelve una 3-tupla $(r, m, X)$, donde $r$ es el nuevo estado de la máquina, $m$ el nuevo símbolo que se ha devuelto en la casilla 
                        determinada, y $X$ es un valor entre L y R, dependiendo si el cabezal debe moverse a izquierda o derecha.
        \item $q_0$ se establece como el estado inicial del que parte la máquina antes de realizar ninguna acción.
        \item $q_A$ es el conjunto de estados de aceptación. Se trata de un subconjunto del conjunto de estados $Q$. En caso de que el programa acabe con un estado de aceptación, se considerará como exitosa la
                        ejecución del programa.
        \item $q_R$ es el conjunto de estados de rechazo, que también pertenece al conjunto de estados $Q$. Al contrario de lo que ocurre con los de aceptación, en caso de que el programe pare su ejecución en 
                        alguno de estos estados, se considerará infructuosa su ejecución.
\end{enumerate}

\section{El funcionamiento}

Una vez entendida la función de transición y todos los componentes de la 7-tupla, planteemos un supuesto de funcionamiento de una máquina de Turing $M$. Como hemos dicho, en una primera instancia
disponemos de una cinta infinita hacia la derecha\footnote{Es decir, finita por la izquierda, con una casilla de inicio. Este supuesto se puede generalizar para una cinta infinita en ambos sentidos, pero en este
ejemplo trataremos con la mencionada.}, donde la máquina recibe como entrada una sucesión de caracteres $\sigma = \sigma_1 \sigma_2 \dots \sigma_n \in \Sigma$ que estarán situados en las $n$ primeras casillas de 
la cinta. Notemos que, como $\Sigma$ no contiene al espacio en blanco, podremos encontrar fácilmente el final de la entrada, al encontrar el primer espacio en la casilla $n+1$.

Teniendo definidas la entrada y todos los componentes de la máquina, podemos proceder con su ejecución. En su proceso, seguirá fielmente todos los pasos y todas las casuísticas definidas en la función de transición:
si encuentras tal símbolo, realiza una determinada acción, desplázate (si es necesario), y en el siguiente paso ejecuta la función de transición para el estado $x$. Y así seguirá hasta que el estado que devuelva
la función sea un estado de parada, ya sea de aceptación o de rechazo.

Habiendo visto un ejemplo de ejecución, en este punto podemos plantearnos ciertas casuísticas peculiares que pueden suceder al aplicar la función de transición definida de una determinada manera. Por ejemplo, en el 
supuesto en el que el cabezal deba moverse hacia la izquierda estando en la primera casilla, podríamos salirnos de la memoria establecida. Este caso se solucionaría de manera sencilla, definiendo que, en caso de que se 
produzca dicha casuística, mantener el cabezal en la misma posición (a pesar de que la función de transición indique lo contrario).

Otro supuesto que se puede dar es que, a priori, nosotros esperamos que en un número finito de pasos la máquina $M$ llegue a algún estado de parada, ya sea de aceptación o de rechazo. Pero, ¿qué pasaría si no alcanzara
ninguno de estos estados en ningún momento de su ejecución? Daría resultado a una máquina que quedaría ejecutándose de manera indefinida\footnote{Véase \cite{sipser2006introduction} para varios ejemplos de máquinas en 
iteraciones finitas e infinitas.}.

¿Podemos a priori saber si nuestra máquina llegará en algún momento a un estado de parada?
Una función de transición puede ser tan compleja como uno quiera definirla. Si tu máquina lleva ejecutándose durante miles de iteraciones, es posible que en el siguiente paso se detenga, o puede ser también que no se detenga
nunca. ¿Hay alguna forma de saberlo (sin tener que ejecutar tu máquina) de si en la máquina está bien definido el criterio de parada?

\section{El problema de la parada}

\subsection{Definiciones previas}

Consideremos nuestra máquina $M$. En cada iteración de su función de transición, se modifica su estado, la información que hay en la cinta y la posición del cabezal. Aunando el cambio de estas tres características, podemos 
definir lo que se conoce como configuración. Una máquina que se ejecuta hasta llegar a un estado de parada en $r$ iteraciones ha recorrido las configuraciones $C_1, \dots, C_r$\footnote{Aquí definimos $C_1$ como configuración 
inicial, con su estado inicial $q_0$, el estado inicial de la cinta y el cabezal en la primera posición de la misma.} hasta detenerse. 

En este sentido, cuando el estado de la configuración $C_r$ sea un estado de aceptación, diremos que $C_r$ es una configuración de aceptación. En caso contrario, cuando el estado de dicha configuración sea de rechazo, 
nos encontraremos ante una configuración de rechazo. Por tanto, diremos que nuestra máquina $M$ \textit{acepta} el estado inicial de la cinta $\sigma$ si existe una sucesión finita de configuraciones de longitud $s$ tales 
que la máquina recorre cada una de ellas de forma secuencial, donde $C_1$ es la configuración inicial con información en la cinta $\sigma$, y $C_s$ es una configuración con un estado de aceptación.\footnote{Podemos
descartar la casuística en la que, para un cierto $i \in \N$ cumpliendo $1 < i < s$, tengamos una configuración $C_1, \dots, C_i, \dots, C_s$ tal que dicha configuración sea de aceptación o de rechazo. En tal caso 
bastará con tomar la subsucesión $C_1, \dots, C_i$ como la sucesión a estudiar si es aceptada o no por la máquina $M$. Por tanto, consideramos que $C_s$ es la primera configuración con un estado de parada de nuestra sucesión.}

Con ello, podríamos considerar para una cierta máquina $M$ el conjunto de todas las entradas $\sigma$ tales que nuestra máquina acepte dicha entrada. A esto lo consideraremos como el lenguaje de $M$, denotado 
por $L(M)$. Teniendo este punto claro, diremos que nuestro lenguaje es \textit{reconocible} por una máquina de Turing $N$ si cualquier elemento del lenguaje es aceptado por $N$\footnote{Podemos notar que, dado un 
lenguaje cualquiera $L_M$, éste será reconocido por la máquina $N$ si y sólo si $L_M \subseteq L(N)$.}.

Como comentábamos anteriormente, muchas veces no es fácil distinguir entre una máquina que, para una cierta entrada, está tardando demasiado en responder, o que simplemente ha entrado en un bucle y correrá de manera indefinida. 
Es cierto que hay máquinas que pueden ser muy útiles a pesar de incurrir en bucles infinitos para ciertas casuísticas, pero a pesar de ello, las máquinas más interesantes a estudiar son las decisoras.

Estas máquinas de Turing tienen la peculiaridad de que, para cualquier entrada perteneciente a cualquier lenguaje, obtienen siempre en un número finito de pasos un estado de parada, ya sea de aceptación o de rechazo.
De la misma manera, existen reducciones de esta definición (reduciendo el conjunto de estudio), que se ajustan a un lenguaje en concreto, dando parada para todos sus elementos. En ese caso, diremos que dicho lenguaje es 
Turing-decidible\footnote{En algunos textos también se los conoce como simplemente decidibles, pero he querido remarcar la existencia de una máquina de Turing que cumpla el criterio.}, siempre y cuando exista al menos 
una máquina que cumpla el criterio mencionado.

\subsection{Algoritmos imposibles}

Llegados a este punto, nos pueden aturdir algunas preguntas naturales. ¿Cómo podemos saber si una máquina es efectivamente decisora? Es decir, ¿existe algún algoritmo capaz de decidir si una máquina (previa a su ejecución)
para una entrada dada se detendrá en un número finito de pasos? Podemos observar cierta recursividad en la propia pregunta. Estamos planteando la existencia de una máquina\footnote{Denotar que, en este caso, recurrimos a la tesis
de Church-Turing para asociar la existencia de algoritmo con la existencia de su respectiva máquina de Turing. Véase \cite{copeland1997church}.} $M_1$ que sea capaz de decidir si otra máquina $M_2$ se detendrá. Es decir, 
una máquina que recibe de entrada el algoritmo de otra máquina. Pero, ¿no tendríamos el mismo problema con la máquina $M_1$? ¿Cómo podemos saber si esta máquina se detendrá?

Este famoso problema es el conocido como problema de la parada (o \textit{halting problem} en inglés). Fue publicado por Turing en su resonado artículo \textit{On Computable Numbers, with an Application to the 
Entscheidungsproblem}\footnote{Véase \cite{turing1936computable}.}\footnote{Recordar al lector Entscheidungsproblem se trata del problema de decibilidad planteado de forma no cerrada por Gödel en sus teoremas de incompletitud, 
donde dejaba pendiente la solución del problema planteado por Hilbert.}, donde consiguió demostrar que el problema de la parada es indecidible, es decir, no existe ninguna máquina $M_1$ cumpliendo las condiciones antes 
exigidas.

Gracias a estos grandes resultados, podemos ver que, para máquinas de Turing relativamente complejas, la etiquetación de una máquina como "decisora" no es tarea fácil, y en algunos casos imposible. Veremos a continuación 
más en profundidad cómo Turing, de una manera sencilla e intuitiva, caminó sobre el problema de manera elegante, dejando huella hasta la actualidad con éste y otros muchos resultados, y plantando la semilla de las bases de 
la computación moderna.



\section{La decibilidad}

En su artículo de 1936, Turing no sólo presentó el problema de la parada, sino que también propuso una aplicabilidad al problema de la decibilidad planteado abiertamente por Hilbert y maquetado por Gödel. Gracias a la 
definición de máquina de Turing y del problema de la parada, se pudo definir la existencia de problemas que no disponen de un proceso válido que lleve a la demostración o refutación del mismo\footnote{En este sentido, 
hacemos referencia de la misma manera al término "algoritmo", apoyándonos en la tesis de Church-Turing.}.

\subsection{Demostración del problema de la parada}

La idea intuitiva detrás de la demostración radica en realizar una demostración por reducción al absurdo. Es decir, suponemos a priori la existencia de una máquina decisora capaz de dar solución al problema de la parada para
cualquier máquina dada. Una vez supuesto esto, bastará con definir una segunda máquina que responda en función de al anterior, y hacer que la máquina decisora interactúe con esta segunda. Llegaremos así a un absurdo y quedará
demostado lo propuesto\footnote{Denotar que las máquinas de Turing, debidamente construidas, son capaces de concatenarse. Es decir, la composición de dos o más máquinas de Turing pueden dar lugar a una nueva máquina, como resultado
de la composición de las anteriores.}.

Comencemos como hemos mencionado. Consideremos una máquina de Turing $P$ cualquiera y una entrada fija $\sigma$, la cual se aplica a la propia máquina $P$ esperando una salida (o una ejecución infinita). Supongamos de la
existencia de una máquina $M$ que, al recibir como entrada a la máquina $P$ con entrada $\sigma$, sea capaz de decidir si dicha máquina $P$ se detendrá en una cantidad finita de iteraciones. Llamemos a esta máquina $M$ como
$DECISORA$ para una mayor sencillez.

Definamos a continuación una nueva máquina $L$ mediante el argumento diagonal de Cantor\footnote{Véase \cite{binder2008theories}.}. Es decir, nuestra máquina $L$ (que llamaremos $DIAGONAL$) es capaz de recibir como entrada 
una máquina cualquiera\footnote{Tengamos en cuenta que siempre que hablemos de máquinas que reciben como entrada otra máquina, hacemos referencia a recibir el algoritmo de una máquina debidamente codificado.}, y realizar una
serie de pasos, que representaremos como una composición de máquinas:

\begin{enumerate}
        \item Dada una entrada $\sigma$, que suponemos como una máquina debidamente codificada, realiza una evaluación de sí misma, dando lugar a un par $\sigma, \sigma$, donde la máquina $\sigma$ recibe como entrada a $\sigma$.
                Este par lo proporciona como salida en la cinta.
        \item Una vez generado el par máquina-entrada, proporcionamos dicho par como entrada en nuestra máquina $DECISORA$, la cual proporcionará como salida un estado de aceptación o de rechazo, dependiendo de si la máquina $\sigma$
                se ha detenido o no. Notar que nuestra máquina $DECISORA$ nunca incurre en bucle.
        \item Por último, consideraremos la salida proporcionada por $DECISORA$, y en caso de que la salida sea un estado de aceptación, nuestra máquina $DIAGONAL$ dará como salida un bucle infinito. En caso contrario, cuando reciba 
                como entrada un estado de rechazo, la máquina proporcionará un estado de aceptación, deteniéndose en un número finito de pasos.
\end{enumerate}

En conjunto, hemos definido la máquina de Turing $DIAGONAL$, la cual recibe como entrada una máquina cualquiera $\sigma$, y proporciona como salida o bien un bucle, o bien un estado de aceptación, siguiendo los pasos anteriormente
descritos.

Notemos que nuestra aplicación $DIAGONAL$, en palabras más sencillas, incurre en el resultado opuesto al que determina $DECISORA$, en una correspondencia biunívoca:

\begin{equation*}
        (DIAGONAL, \sigma) \textit{ entra en estado de aceptación } \Leftrightarrow (DECISORA, \sigma) \textit{ entra en estado de rechazo.}\footnote{Se entiende por el par
        $(M, \sigma)$ a la evaluación de la máquina $M$ para una entrada $\sigma$.}
\end{equation*}

O dicho en otras palabras, y continuando con la equivalencia, 

\begin{equation*}
        (DIAGONAL, \sigma) \textit{ se detiene } \Leftrightarrow (\sigma, \sigma) \textit{ no se detiene.}
\end{equation*}

Y esto ocurre para cualquier máquina $\sigma$ que nosotros consideremos. Tomemos por tanto $\sigma = DIAGONAL$ (esto lo podemos hacer ya que $DIAGONAL$ es una máquina de Turing, por ser composición de éstas). Por tanto, sustituyendo en el 
resultado anterior tenemos lo siguiente:

\begin{equation*}
        (DIAGONAL, DIAGONAL) \textit{ se detiene } \Leftrightarrow (DIAGONAL, DIAGONAL) \textit{ no se detiene.}
\end{equation*}

Hemos llegado a una contradicción, y ello se debe a la única suposición que hemos hecho en todo nuestro proceso lógico: la existencia de la máquina $DECISORA$. Por tanto, por reducción al absurdo, hemos demostrado la no existencia 
máquinas con estas características.

%\subsection{Otros problemas indecidibles}


\chapter{El estado del programa de Hilbert}

Del problema original quedan pocos pilares en pie. Tras el desmoronamiento de la completitud y de la consistencia, Turing consiguió 
con un concepto tan simple y tan revolucionario al mismo tiempo, darle el golpe de gracia a la única piedra que quedaba en pie.

Cierto es que con otros muchos problemas del programa original de Hilbert no ha sucedido lo mismo, pero por suerte o por desgracia
este problema no tuvo el fin que Hilbert esperaba. Ya en un inicio, los resultados de Gödel habían dado un refresco al concepto propio
de lo que conocemos como matemáticas. Había revolucionado la forma de estudiar los problemas, y de poder tener clara la incompletitud
de los sistemas formales sobre los que trabajamos. Puede ser que por esta parte Hilbert no viera resolución a su propuesta, pero a día 
de hoy podemos agradecer a Gödel su grandes aportaciones a la matemática moderna. 

A pesar de ello, Hilbert mantuvo hasta el final su fe en la decibilidad, pero tampoco consiguió llegar el mismo a buen puerto. Turing, 
con un concepto simple, novedoso y revolucionario, demostró la existencia de problemas indecidibles, y por tanto la posibilidad de existencia
de problemas que quedarán irresolubles de por vida. Puede parecer una desgracia para Hilbert, pero probablemente haya sido de los mayores 
logros para la matemática y la computación moderna. Gracias a este resultado, conocemos el mundo tal y como lo conocemos. 

A fin de cuentas, el problema de Hilbert ha sido resuelto, aunque no haya sido en la vía que Hilbert esperaba. El lado bueno de este 
tipo de resultados, es que la no existencia de solución esperada no siempre significa un fracaso. Tenemos aquí, en la misma línea temporal
de un mismo problema, dos ejemplos claros de problemas ``sin solución'' que han supuesto un cambio generacional en las matemáticas. Si hubiésemos demostrado lo contrario, nos encontraríamos en un estado del arte distinto, en un mundo con problemas diferentes que resolver,
y aportando soluciones desde perspetivas para nosotros desconocidas. Pero no ha sucedido así. La búsqueda de soluciones alternativas a los 
problemas nos han hecho descubrir contraejemplos que nos han hecho llegar hasta donde hemos llegado, tanto como sociedad humana como matemática.

Ambas aportaciones, tanto las de Turing como las de Gödel, nos han ayudado a descubrir la existencia de matemáticas perfectamente imperfectas,
con problemas tan difíciles que nos hacen cuestionarnos si en realidad tendrán solución. Gracias a Hilbert y a su propuesta del problema, 
hemos descubierto cosas que de ninguna otra forma hubiéramos descubierto, dando pasos hacia adelante hacia la construcción de una matemática
completamente incompleta, perfectamente imperfecta y con la magia de la duda.

\endinput
%------------------------------------------------------------------------------------
% FIN DEL CAPÍTULO. 
%------------------------------------------------------------------------------------

